% !PRINTED 249
\weberpage{249}

Finally, there remains the condition for the equality of all four roots
\weq{13}{a = 0, \quad b = 0, \quad c = 0.}

These relations can be made very vivid by a geometric interpretation, and although geometric considerations do not properly lie within our plan, we shall nevertheless not neglect to point them out to the reader from time to time.

If one interprets $a$, $b$, $c$ as rectangular coordinates in space, then every point of space is to be regarded as the carrier of a certain biquadratic equation of the form (1), $f(x) = 0$; all equations that have a given number $x$ as a root are represented by the points of a plane $[f(x) = 0]$. The equation $D = 0$ is the equation of a curved surface (of the fifth degree), the discriminant surface, which is generated by the lines of intersection of the planes $f(x) = 0$, $f'(x) = 0$, and is therefore a developable surface. It is the envelope of all the planes $f(x) = 0$.

The surface has an edge of regression consisting of two branches, which is represented by the equations (11), (12), and a double line, which is determined by the equations $b = 0$, $a^2 - 4c = 0$ and therefore has the shape of a parabola. On the side of negative $a$ the surface crosses itself in this parabola. On the side of positive $a$ the parabola continues as an isolated line.

The discriminant surface divides the whole of space into three compartments, of which two are connected only along the double parabola, and these compartments contain the points to which no, two, and four real roots correspond. If for the moment we call these three compartments $[0]$, $[2]$, $[4]$, then $[0]$ borders $[4]$ only along the double parabola, while $[0]$ as well as $[4]$ border $[2]$ along the parts of the surface.

% [sic] The source defines [4,2] but here prints the second boundary pair as [0],[4]; preserved diplomatically.
We shall denote by $[0, 2]$ and $[4, 2]$ the boundary surfaces of $[0]$, $[2]$ and of $[0]$, $[4]$. The isolated part of the parabola continues into the interior of $[0]$. The edges of regression lie on the part of the surface that separates $[4]$ from $[2]$.

At the origin of coordinates the three parts of space and their boundary surfaces and boundary lines meet. The points of the surface parts $[0, 2]$ represent equations with two equal and
%CONT>>
